% Why a linear seasonal baseline is accurate at every bucketing granularity.
% Compile: pdflatex granularity_invariance.tex
\documentclass[11pt]{article}
\usepackage[margin=1in]{geometry}
\usepackage{amsmath,amssymb,amsthm}
\usepackage{xcolor}
\usepackage[colorlinks=true,linkcolor=blue,urlcolor=blue]{hyperref}
\setlength{\emergencystretch}{3em}

\theoremstyle{plain}
\newtheorem{proposition}{Proposition}
\newtheorem{corollary}{Corollary}
\theoremstyle{remark}
\newtheorem{remark}{Remark}

\newcommand{\R}{\mathbb{R}}
\newcommand{\E}{\mathbb{E}}
\newcommand{\Var}{\operatorname{Var}}
\newcommand{\Cov}{\operatorname{Cov}}

\title{\textbf{Why a Linear Seasonal Baseline is Accurate\\ at Every Bucketing Granularity}}
\author{Traffic Forecasting Project}
\date{June 2026}

\begin{document}
\maketitle

\begin{abstract}
\noindent
The calendar-adjusted seasonal baseline -- a recency-weighted average of the same
hour-of-week, times a calendar factor -- is empirically accurate whether applied to
a single sensor at hourly resolution, a region total, or a weekly aggregate. We
explain why with two short results. \emph{Equivariance:} the baseline is a linear
operator and therefore commutes with any linear aggregation, so the same model is
automatically self-consistent at every granularity (forecasting then aggregating
equals aggregating then forecasting). \emph{Signal-to-noise scaling:} the
forecastable part is a low-rank, shared seasonal signal while the residual is
high-rank idiosyncratic noise, so aggregation amplifies signal coherently
($\propto n$) and averages noise incoherently ($\propto \sqrt{n}$), making the same
model at least as accurate -- usually more -- as the buckets coarsen.
\end{abstract}

\section{Setup}

Let $y_{i,t}$ be the series for unit $i\in\{1,\dots,N\}$ (e.g.\ a sensor) at hour
$t$, with weekly period $P=168$. The \emph{seasonal baseline} is the linear filter
\begin{equation}
\hat y_{i,t} \;=\; (B y)_{i,t} \;=\; \sum_{k=1}^{K} w_k\, y_{i,\,t-Pk},
\qquad w_k\ge 0,\quad \sum_{k=1}^{K} w_k = 1,
\label{eq:B}
\end{equation}
i.e.\ a (possibly decaying) average over the same hour-of-week in the previous $K$
weeks. The flat 4-week rule is $K=4$, $w_k=\tfrac14$; the EWMA rule is
$w_k\propto\alpha^{k-1}$. The crucial structural facts about $B$ are that it is
\textbf{linear}, \textbf{time-invariant}, and uses \textbf{the same weights for
every unit $i$}.

A \emph{granularity} is produced by a linear \emph{aggregation operator} $A$ that
maps the fine series to a coarser one by summing over blocks of indices:
\begin{align}
\text{spatial:}&\quad (A^{\mathrm{sp}} y)_{R,t} = \sum_{i\in R} y_{i,t}
  \quad(\text{sum over a set of units }R),\\
\text{temporal:}&\quad (A^{\mathrm{tw}} y)_{i,w} = \sum_{h=0}^{P-1} y_{i,\,Pw+h}
  \quad(\text{sum the }168\text{ hours of week }w).
\end{align}
(Means are the same up to a constant, so the argument is identical for averages.)

We decompose the series into a predictable seasonal/level signal and zero-mean
noise,
\begin{equation}
y_{i,t} \;=\; \mu_{i,t} \;+\; \varepsilon_{i,t},
\qquad \mu_{i,t}=\ell_{i,t}\, s_{i,\phi(t)},\quad \E[\varepsilon_{i,t}]=0,
\label{eq:decomp}
\end{equation}
where $\phi(t)=t\bmod P$ is the hour-of-week, $s_{i,\cdot}$ is a (slowly varying)
weekly \emph{shape}, $\ell_{i,t}$ a slow \emph{level}, and $\varepsilon$ is the
idiosyncratic, incident-driven residual. The baseline $B$ is a near-unbiased,
low-variance estimator of $\mu$ (it averages recent same-phase values); its error
is dominated by the irreducible noise $\varepsilon$.

\section{Equivariance: one model at every granularity}

\begin{proposition}[Aggregation equivariance]\label{prop:equi}
For spatial aggregation, $B$ commutes with $A^{\mathrm{sp}}$ exactly:
\[
B\,A^{\mathrm{sp}} \;=\; A^{\mathrm{sp}}\,B .
\]
For whole-week temporal aggregation, the hourly baseline summed over a week equals
the \emph{same-form} weekly baseline applied to the weekly totals:
\[
A^{\mathrm{tw}} B \;=\; B'\, A^{\mathrm{tw}},
\]
where $B'$ has identical weights $w_k$ but at weekly lags.
\end{proposition}

\begin{proof}
Spatial, using linearity and that $w_k$ does not depend on $i$:
\begin{align*}
(B A^{\mathrm{sp}} y)_{R,t}
&= \sum_{k} w_k (A^{\mathrm{sp}} y)_{R,t-Pk}
 = \sum_{k} w_k \sum_{i\in R} y_{i,t-Pk}\\
&= \sum_{i\in R} \sum_{k} w_k\, y_{i,t-Pk}
 = (A^{\mathrm{sp}} B y)_{R,t}.
\end{align*}
Temporal, using that the lags $Pk$ are whole multiples of the aggregation block $P$:
\begin{align*}
(A^{\mathrm{tw}} B y)_{i,w}
&= \sum_{h=0}^{P-1} \sum_{k} w_k\, y_{i,\,Pw+h-Pk}
 = \sum_{k} w_k \sum_{h=0}^{P-1} y_{i,\,P(w-k)+h}\\
&= \sum_{k} w_k\, (A^{\mathrm{tw}} y)_{i,\,w-k}
 = (B' A^{\mathrm{tw}} y)_{i,w}. \qquad\qedhere
\end{align*}
\end{proof}

\begin{corollary}[Self-coherence]\label{cor:coh}
The fine-level forecasts sum exactly to the coarse-level forecast:
$\sum_{i\in R}\hat y_{i,t} = \widehat{(A^{\mathrm{sp}}y)}_{R,t}$, and likewise for
time. Hence \emph{a single model is simultaneously the model at every aggregation
level}: one need not refit a different form per granularity, and forecasts made at
different resolutions are automatically consistent.
\end{corollary}

\begin{remark}[The multiplicative calendar factor]
The calendar-adjusted forecast is $\hat y^{\mathrm{adj}}_{i,t}=f_{i,t}\,(By)_{i,t}$.
Aggregating, $\sum_{i\in R} f_{i,t}(By)_{i,t}=f_{R,t}\sum_{i\in R}(By)_{i,t}$ holds
\emph{exactly} when the factor is uniform across the aggregated units,
$f_{i,t}\equiv f_{R,t}$. A \emph{pooled} factor (one value per
(holiday,\,hour-of-day), shared across units) satisfies this, so the calendar
correction is coherent too; a heterogeneous per-unit factor makes the aggregate
factor a baseline-weighted average of the unit factors, so equivariance becomes
approximate. The linear backbone \eqref{eq:B} is always exactly coherent.
\end{remark}

\begin{remark}[Same weights]
Equivariance is exact only if the \emph{same} weights $w_k$ are used at each level.
Re-tuning the decay per granularity breaks it slightly; in practice a shared
$\alpha$ keeps the hierarchy coherent.
\end{remark}

\section{Signal-to-noise improves under aggregation}

Equivariance says the same model \emph{form} is correct everywhere; it does not by
itself say the model is \emph{accurate} everywhere. Accuracy follows from how signal
and noise scale under aggregation.

Write the aggregate of $n$ units (a region, or one unit's weekly total over $n=P$
hours) and split it via \eqref{eq:decomp} into an aggregate signal and an aggregate
noise. Two facts distinguish them:
\begin{itemize}
\item \textbf{The signal is shared (coherent).} Every unit carries the \emph{same}
  weekly rhythm $s_{\phi(t)}$, so the seasonal components add in phase: the aggregate
  signal grows like $n\,\bar\mu$, and its variation across $t$ grows like $n$ times
  the per-unit variation.
\item \textbf{The noise is idiosyncratic (incoherent).} Incidents are local. With
  per-unit noise variance $\sigma^2$ and pairwise correlation $\rho$,
  \begin{equation}
  \Var\!\Big(\textstyle\sum_{i=1}^{n}\varepsilon_{i,t}\Big)
  = n\sigma^2 + n(n-1)\rho\,\sigma^2
  = n\sigma^2\big[1+(n-1)\rho\big].
  \label{eq:noisevar}
  \end{equation}
\end{itemize}

\begin{proposition}[SNR scaling]\label{prop:snr}
Let the aggregate signal scale as $n$ and the aggregate noise standard deviation be
$\sigma\sqrt{n[1+(n-1)\rho]}$ from \eqref{eq:noisevar}. Then the
\emph{relative} forecast error (noise-to-signal) of the baseline obeys
\[
\frac{\text{aggregate noise s.d.}}{\text{aggregate signal}}
\;\sim\; \frac{\sigma}{\bar\mu}\,\sqrt{\frac{1+(n-1)\rho}{n}}
\;\xrightarrow[n\to\infty]{}\;
\begin{cases}
0, & \rho = 0 \ (\text{idiosyncratic noise}),\\[4pt]
\dfrac{\sigma}{\bar\mu}\sqrt{\rho}, & \rho > 0 \ (\text{common component}).
\end{cases}
\]
Equivalently, the unexplained-variance fraction $1-R^2$ shrinks monotonically as the
buckets coarsen, to $0$ if the residual noise is purely idiosyncratic and to a
floor $\propto\rho$ if it has a shared component.
\end{proposition}

\begin{proof}
Divide the aggregate noise s.d.\ $\sigma\sqrt{n[1+(n-1)\rho]}$ by the aggregate
signal $\sim n\bar\mu$ to get $\frac{\sigma}{\bar\mu}\sqrt{(1+(n-1)\rho)/n}$. For
$\rho=0$ this is $\frac{\sigma}{\bar\mu}n^{-1/2}\to0$; for $\rho>0$,
$(1+(n-1)\rho)/n\to\rho$, giving the stated floor. Monotonicity in $n$ follows since
$(1+(n-1)\rho)/n$ is decreasing for $\rho<1$.
\end{proof}

\noindent
So coarser granularity carries a higher signal-to-noise ratio: the same seasonal
model -- which captures $\mu$ and leaves $\varepsilon$ -- explains a larger fraction
of variance. (Empirically the weekly $R^2$ exceeded the hourly $R^2$; part of that
is this genuine noise-averaging, and part is the larger seasonal variance inflating
the $R^2$ denominator -- the former is the real effect.)

\section{Synthesis}

The two results compose. Using the factor structure $\mu_{i,t}=\ell_{i,t}s_{i,\phi(t)}$:
\begin{itemize}
\item the signal $\mu$ is \textbf{low-rank and shared} -- a common deterministic
  seasonal pattern, present at \emph{every} scale (self-similar: a weekly cycle
  within a week, an annual cycle across weekly totals);
\item the noise $\varepsilon$ is \textbf{high-rank and idiosyncratic}.
\end{itemize}
The baseline $B$ is a near-unbiased, low-variance estimator of $\mu$ at any scale.
Aggregation then does two complementary things: by Proposition~\ref{prop:equi} it
\emph{preserves the model's form and coherence}, and by
Proposition~\ref{prop:snr} it \emph{amplifies the shared signal while averaging out
the idiosyncratic noise}. Equivariance gives consistency across granularities; SNR
scaling gives (weakly) increasing accuracy. Together:
\begin{center}
\fbox{\begin{minipage}{0.9\linewidth}\centering\vspace{2pt}
\emph{one simple, correctly-specified linear seasonal model is automatically
self-consistent and accurate at every bucketing granularity.}
\vspace{2pt}\end{minipage}}
\end{center}

\paragraph{Caveats, all orthogonal to granularity.} Exact equivariance needs the
same weights at each level and (for the multiplicative factor) a pooled, uniform
calendar factor. Correlated residual noise ($\rho>0$) sets a floor on the
SNR gain. And the \emph{bias} properties are scale-independent: the baseline's level
estimate goes stale as the forecast horizon grows, and its structural assumption
breaks at regime shifts (e.g.\ a demand shock) -- these inflate error at \emph{all}
granularities, not at one.

\end{document}
